\documentclass[../RFC-HDFG-2026-003.tex]{subfiles}

\begin{document}

\begin{abstract}


This RFC proposes in-file blob configuration storage for HDF5 I/O
filters: a binary channel for filter configuration data too large for
the \texttt{cd\_values} array or the parameter-string API introduced by
RFC-HDFG-2026-001. Two use cases motivate it: a JIT-compiled filter
module that must store multi-megabyte pre-processed source code
alongside the dataset, and a filter that needs only a small reference
to a separately-stored HDF5 dataset (e.g., a spatial mask).

The proposed design adds three new, NULL-able callbacks to
\texttt{H5Z\_class3\_t} --- \texttt{write\_blob}, \texttt{read\_blob}, and
\texttt{close\_blob} --- and a new public API, \texttt{H5Pappend\_filter\_blob},
that appends a filter to a pipeline together with an opaque byte buffer.
The blob is persisted as a global-heap (\texttt{H5HG}) object referenced
from a new pipeline-message format version; when a filter does not
implement custom storage callbacks, the library falls back to a default
\texttt{H5HG}-based reader/writer. This RFC resolves the parallel I/O
write protocol, the \texttt{H5Pencode}/\texttt{H5Pdecode} serialization
policy, \texttt{h5repack} cross-file copy behavior, and the on-disk
storage format, and shows that both motivating use cases are covered by
the same callback signatures without additional public API surface.
\end{abstract}

\vfill
{\footnotesize\textbf{Acknowledgements:} This work is supported by the U.S.\
Department of Energy, Office of Science, Advanced Scientific Computing Research,
under Contract DE-AC02-06CH11357 and by the U.S.\ Department of Energy, Office
of Science, Office of Advanced Scientific Computing Research, Scientific
Discovery through Advanced Computing (SciDAC) program.}

\end{document}
